\documentclass[report.tex]{subfiles}

\begin{document}

\section{The Free Simulation}

While explicit simulations effectively prevent infinite stuttering,
they impose a substantial proof burden.
Every stuttering transition requires exhibiting a strictly decreasing
measure in a well-founded order.
Constructing such a measure can be difficult in practice,
particularly when compiler optimizations rearrange loop structures.
In such scenarios,
finding a global metric that strictly decreases on every asynchronous step
while remaining invariant across complex control transfers
is often impractical.

To overcome this rigidity,
we adapt the \emph{Stuttering for Free} framework
introduced by \textcite{cho-23-stuttering}.
We extend their foundational technique to accommodate
skewed simulation relations.
All underlying theoretical results have been formalized and verified
within our Rocq development.

\subsection{Implicit Simulation}

An effective approach to eliminating explicit stuttering measures
is to combine inductive and coinductive reasoning within a single relation.
In this formulation,
finite sequences of asynchronous execution steps are justified inductively,
whereas mutual progress steps are justified coinductively.
Because explicit stuttering indices are entirely omitted from the relation,
this family of relations is referred to as \emph{implicit simulations}
by \textcite{cho-23-stuttering}.
We define an implicit simulation relation,
denoted $\isim{t}{s}{\Phi}$,
whose inference rules are presented in \cref{fig:implicit-sim}.

\begin{figure}[htb]
  \begin{align*}
    &
    \begin{prooftree}
      \hypo{\bfinal{\Phi}{t}{s}}
      \infer1[\textsc{IRelated}]{\isim{t}{s}{\Phi}}
    \end{prooftree}
    &
    &
    \begin{prooftree}
      \hypo{\stuck[s]{s}}
      \infer1[\textsc{ISourceStuck}]{\isim{t}{s}{\Phi}}
    \end{prooftree}
  \end{align*}
  \begin{equation*}
    \begin{prooftree}
      \hypo{\step[s]{s}{s'}}
      \hypo{\isim{t}{s'}{\Phi}}
      \infer2[\textsc{ISourceSteps}]{\isim{t}{s}{\Phi}}
    \end{prooftree}
  \end{equation*}
  \begin{equation*}
    \begin{prooftree}
      \hypo{\canp[t]{t}}
      \hypo{\forall t', \step[t]{t}{t'} \implies \isim{t'}{s}{\Phi}}
      \infer2[\textsc{ITargetSteps}]{\isim{t}{s}{\Phi}}
    \end{prooftree}
  \end{equation*}
  \begin{equation*}
    \begin{prooftree}
      \hypo{\canp[t]{t}}
      \hypo{\forall t',
        \step[t]{t}{t'} \implies
        \exists s', \step[s]{s}{s'} \land \isim{t'}{s'}{\Phi}}
      \infer[coind]2[\textsc{IBothSteps}]{\isim{t}{s}{\Phi}}
    \end{prooftree}
  \end{equation*}
  \caption{Derivation rules of the \emph{skewed implicit simulation}.}
  \label{fig:implicit-sim}
\end{figure}

The rule \textsc{IRelated} characterizes terminating executions,
mirroring the \textsc{ERelated} rule of the explicit simulation
presented in \cref{fig:explicit-sim}.
Stuttering is soundly handled
without requiring an explicit decreasing measure.
Because \textsc{ISourceSteps} and \textsc{ITargetSteps} are inductive
constructors,
they permit only finitely many consecutive asynchronous transitions.
The coinductive constructor \textsc{IBothSteps} requires
both programs to execute a synchronized step,
removing the inductive proof obligation
and enabling unbounded mutual progress.

\subsection{Step-Indexed Free Simulation}

While the implicit simulation eliminates explicit stuttering measures,
it remains restrictive:
the coinductive rule \textsc{IBothSteps}
demands that both programs take a step at the exact moment
when coinductive reasoning is resumed.
Removing this lock-step synchronization requirement
leads to simpler invariants
and more direct verification proofs~\parencite{cho-23-stuttering}.
This decoupling is the core idea
developed in the \emph{Stuttering for Free} framework.

The \emph{free simulation},
denoted $\fsim{t}{j}{i}{s}{\Phi}$,
overcomes this limitation by decoupling execution steps
from coinductive progress.
The relation is parameterized by two well-founded index domains,
$(J, \prec_J)$ for the target side
and $(I, \prec_I)$ for the source side.
These domains track the progress accumulated by each side
since the previous coinductive step.
The formal inference rules of the skewed free simulation
are presented in \cref{fig:free-sim}.
Like the implicit simulation,
it accommodates arbitrary postconditions $\Phi$,
and uses the rule \textsc{FSourceStuck}
to exploit any undefined behavior in the source program.

\begin{figure}[htb]
  \begin{align*}
    &
    \begin{prooftree}
      \hypo{\bfinal{\Phi}{t}{s}}
      \infer1[\textsc{FRelated}]{\fsim{t}{j}{i}{s}{\Phi}}
    \end{prooftree}
    &
    &
    \begin{prooftree}
      \hypo{\stuck[s]{s}}
      \infer1[\textsc{FSourceStuck}]{\fsim{t}{j}{i}{s}{\Phi}}
    \end{prooftree}
  \end{align*}
  \begin{equation*}
    \begin{prooftree}
      \hypo{\step[s]{s}{s'}}
      \hypo{\fsim{t}{j}{i'}{s'}{\Phi}}
      \infer2[\textsc{FSourceSteps}]{\fsim{t}{j}{i}{s}{\Phi}}
    \end{prooftree}
  \end{equation*}
  \begin{equation*}
    \begin{prooftree}
      \hypo{\canp[t]{t}}
      \hypo{\forall t',
        \step[t]{t}{t'} \implies
        \exists j', \fsim{t'}{j'}{i}{s}{\Phi}}
      \infer2[\textsc{FTargetSteps}]{\fsim{t}{j}{i}{s}{\Phi}}
    \end{prooftree}
  \end{equation*}
  \begin{equation*}
    \begin{prooftree}
      \hypo{j' \prec_J j}
      \hypo{i' \prec_I i}
      \hypo{\fsim{t}{j'}{i'}{s}{\Phi}}
      \infer[coind]3[\textsc{FProgress}]{\fsim{t}{j}{i}{s}{\Phi}}
    \end{prooftree}
  \end{equation*}
  \caption{Derivation rules of the \emph{skewed free simulation}.}
  \label{fig:free-sim}
\end{figure}

The design of the free simulation
relies on a fundamental shift in perspective regarding indices.
In an explicit simulation,
an index represents an upper bound on \emph{future} stuttering steps.
In the free simulation,
indices instead record \emph{past} progress
accumulated since the previous coinductive step.
Consequently,
a simulation statement with smaller indices
is stronger than one with larger indices.
Indices are meaningful primarily
during the coinductive proof of the simulation.
Once the simulation is established,
it holds independently of the chosen indices,
a property we justify formally later in \cref{thm:index-irrel}.

The source and target programs take transitions asynchronously
via the inductive rules \textsc{FTargetSteps} and \textsc{FSourceSteps}.
Whenever a program takes a step,
its corresponding step index is updated ($j'$ or $i'$).
Importantly,
$j'$ and $i'$ are not required to be
immediate successors of $j$ and $i$;
instead,
one is free to select any new index,
whether smaller or larger.

This flexibility does not compromise soundness.
The only rule that enables coinductive reasoning is \textsc{FProgress},
which requires both indices to strictly decrease,
from $j$ down to $j' \prec_J j$
and from $i$ down to $i' \prec_I i$.
Because the relations $\prec_J$ and $\prec_I$ are well-founded,
establishing a simulation between non-terminating executions
demands that \textsc{FProgress} be applied infinitely often.
This,
in turn,
requires indices to be repeatedly increased
through \textsc{FTargetSteps} and \textsc{FSourceSteps}.
This structure decouples stepping from coinductive synchronization,
allowing multiple asynchronous transitions
on either side to be matched flexibly.

Moreover,
coinductive reasoning over the free simulation
is governed by the following principle:
\begin{lemma}[Coinduction Principle for Free Simulation]
  \label{lem:fsim-coind}
  To prove that $\fsim{t}{j}{i}{s}{\Phi}$,
  one may assume the coinductive hypothesis for strictly larger indices:
  \begin{equation*}
    \forall j' \succ_J j,\; \forall i' \succ_I i,\quad
    \fsim{t}{j'}{i'}{s}{\Phi}.
  \end{equation*}
\end{lemma}
Intuitively,
this principle allows one to establish
that two programs simulate each other
by assuming that the simulation relation already holds
for the current states whenever both sides have made strict progress.
A concrete example illustrating step-indexed reasoning
in the free simulation
is presented in \cref{app:free-sim-example}.

\subsection{Properties}

As observed earlier,
proving a free simulation for smaller indices
yields a stronger statement
that implies the simulation for larger indices.
This monotonicity property is useful during coinductive proofs.
Formally:
\begin{lemma}[Monotonicity of Free Simulation]
  \label{lem:fsim-mono}
  If $\fsim{t}{j}{i}{s}{\Phi}$ holds,
  then for all indices $j \prec_J j'$ and $i \prec_I i'$,
  and any weaker postcondition $\Phi \implies \Phi'$,
  the simulation $\fsim{t}{j'}{i'}{s}{\Phi'}$ holds.
\end{lemma}

The free simulation is sound
with respect to skewed behavioral refinement:
\begin{theorem}[Soundness of Free Simulation]
  \label{thm:fsim-soundness}
  Let $P_t$ and $P_s$ be programs,
  with initial states $t$ and $s$,
  respectively.
  If $\fsim{t}{j}{i}{s}{=}$ holds for some indices $j$ and $i$,
  then $P_t \slantsqsubseteq P_s$.
\end{theorem}

In addition to soundness,
we established the formal equivalence
among the three simulation frameworks presented:
explicit simulation (\cref{fig:explicit-sim}),
implicit simulation (\cref{fig:implicit-sim}),
and free simulation (\cref{fig:free-sim}).
This equivalence is established through a chain of cyclic implications:
\begin{equation*}
  \text{Explicit Simulation}
  \overset{\text{\ding{172}}}{\implies}
  \text{Implicit Simulation}
  \overset{\text{\ding{173}}}{\implies}
  \text{Free Simulation}
  \overset{\text{\ding{174}}}{\implies}
  \text{Explicit Simulation}
\end{equation*}
The implications \ding{172} and \ding{173}
proceed directly by coinduction.
The final implication \ding{174} is the most intricate:
it requires synthesizing a well-founded stuttering measure
from the derivation tree of the free simulation.
A direct consequence of this equivalence is the index irrelevance theorem:
\begin{theorem}[Index Irrelevance]
  \label{thm:index-irrel}
  Let $(J, \prec_J)$ and $(I, \prec_I)$
  be well-founded relations without isolated elements%
  \footnote{An element is isolated
    if it has no predecessors or successors in the relation.}.
  If there exist indices $j \in J$ and $i \in I$
  such that $\fsim{t}{j}{i}{s}{\Phi}$,
  then for any other well-founded structures
  $(J', \prec_{J'})$ and $(I', \prec_{I'})$,
  and any indices $j' \in J'$ and $i' \in I'$,
  the relation $\fsim{t}{j'}{i'}{s}{\Phi}$ holds.
\end{theorem}

This theorem cannot be applied within a proof by coinduction%
\footnote{
  Formally,
  when free simulation is characterized as a greatest fixed point,
  monotonicity (\cref{lem:fsim-mono}) holds
  for all intermediate post-fixpoints of the coinductive proofs.
  In contrast,
  index irrelevance (\cref{thm:index-irrel}) is a property
  of the greatest fixed point itself,
  and thus cannot be assumed as a hypothesis
  within an unclosed coinductive derivation.
}.
This theorem highlights a fundamental aspect
of step indices in the free simulation:
they carry no semantic weight once the simulation relation is established.
Their role is purely to ease the proof of the simulation.
They are used to synchronize transitions and bound asynchronous stuttering.

Our statement refines the index irrelevance theorem
formulated by \textcite{cho-23-stuttering}.
Their original formulation,
phrased in terms of ordinals,
implies that well-founded structures
contain no isolated elements.
However,
isolated elements could cause \textsc{FProgress}
transitions to be uninhabited,
preventing the equivalence.
This requirement was identified
during our Rocq formalization of the metatheory of the free simulation.

This theorem implies that any free simulation derivation
can be carried out using the two-element boolean order
$\mathbb{B} = \{\false \prec \true\}$:
\begin{equation*}
  (\exists j \in J, i \in I,\; \fsim{t}{j}{i}{s}{\Phi})
  \iff
  (\exists j, i \in \mathbb{B},\; \fsim{t}{j}{i}{s}{\Phi}).
\end{equation*}
Under the boolean order,
$\false$ indicates that a program has not yet taken a step,
whereas $\true$ signals that progress has occurred.
In principle,
this shows that free simulation proofs
do not require tracking exact step counts,
but only whether progress has occurred
since the previous coinductive step.
Nevertheless,
in the remainder of our work,
we adopt the domain $(\mathbb{N}, <)$
as our index domain.
Indeed,
indices in $\mathbb{N}$ correspond directly
to the number of steps executed during symbolic evaluation,
providing clear intuition that facilitates
interactive verification in Rocq.

\end{document}

% LocalWords:  Coinductive coinductive asymmetricity arities
